% <<<PDFPAGE 21>>>
\chapter{推测解码}

\section{起草与验证}

推测解码（Leviathan 等，2023，以及 Chen 等，2023）利用了一种不对称性。要\emph{验证}大模型是否会生成给定的词元序列，只需一次并行的前向传播。若要让大模型以自回归方式\emph{生成}这段序列，则每生成一个词元都要做一次前向传播。

算法的流程很简洁。
\begin{enumerate}
  \item 用一个小型\textbf{草稿模型}（draft model，30 至 70 亿参数）以较低成本自回归地生成 \(K\) 个候选草稿词元。
  \item 用大型\textbf{目标模型}（target model）做一次前向传播，并行处理这 \(K\) 个草稿词元。
  \item 采用拒绝采样（rejection sampling），依据目标模型的分布\textbf{验证}每个草稿词元。接受符合条件的词元，遇到拒绝时执行校正步骤。
  \item 可以证明，所得结果与直接从目标模型的分布中采样等价。
\end{enumerate}

期望加速比为 \(1/(1-\alpha^K)\)，其中 \(\alpha\) 为单个词元的接受率（acceptance rate）。当 \(\alpha=0.8\)、\(K=5\) 时，每步接受的词元数期望约为 \(3.3\)，因此理论上的解码吞吐量为原来的 \(3.3\) 倍。

\translatornote{正文的原式代入该组数值约为 \(1.49\)，与例中的 \(3.3\) 不符。若每个草稿词元在前面的词元均被接受时的条件接受概率恒为 \(\alpha\)，每轮输出词元数的期望为附录 A 所列的 \((1-\alpha^{K+1})/(1-\alpha)\)。这包含首次拒绝时采样的校正词元，或全部接受时额外采样的一个词元。代入本例约为 \(3.69\)，其中约 \(2.69\) 个是接受的草稿词元。实际加速比还取决于草稿生成和验证的耗时。见 \href{https://arxiv.org/abs/2211.17192}{Leviathan 等（2023）第 3 节}。}

\section{词元树验证}

基本的推测解码只生成一条草稿序列。\textbf{基于树的推测解码}（tree-based speculative decoding）生成一棵由草稿序列组成的\emph{树}。在每个位置，草稿模型提出多个候选词元，形成分支。目标模型使用树形掩码（tree mask）控制注意力，在一次前向传播中验证所有分支。这样同时探索多个候选方案，可显著提高获得较高接受率的概率。

\section{Medusa 多头推测解码}

Medusa（Cai 等，2024）在目标模型上增加 \(K\) 个“Medusa 头”（Medusa heads），每个头都是轻量的线性层，因而省去独立的草稿模型。第 \(k\) 个头在同一次前向传播中预测位置 \(t+k\) 的词元，也就是往后第 \(k\) 个词元。
\[
  \operatorname{logits}_{t+k}=\operatorname{MedusaHead}_k(\mathbf{h}_t)
\]
其中 \(\mathbf{h}_t\) 为位置 \(t\) 的最后一层隐状态。这使模型能自行提出多个后续词元作为草稿。Medusa-2 加入自蒸馏（self-distillation），进一步改善草稿质量，在生成基准测试中达到 \(2\) 至 \(3\) 倍的速度。

% <<<PDFPAGE 22>>>
\section{EAGLE 与高效自回归生成}

Li 等（2024）提出的提高语言模型效率的外推算法（Extrapolation Algorithm for Greater Language-model Efficiency，EAGLE）针对 Medusa 的一个重要局限进行了改进。只根据最后一个隐状态预测后续词元，会丢失关于具体词元序列的信息，从而限制接受率。EAGLE 训练一个轻量草稿模型，同时接收隐状态\emph{和}已生成词元的特征嵌入，因此能更准确地预测下一个词元。EAGLE-2 又引入\textbf{动态草稿树}（dynamic draft tree），根据草稿预测的置信度调整树的结构，在编程和推理任务上达到原来 \(3\) 至 \(4\) 倍的速度。

\section{推测解码的适用条件与局限}

\begin{conceptbox}{推测解码发挥作用的条件}
\begin{itemize}
  \item \textbf{单流、低时延服务}。批大小为 1 时，目标模型未能充分利用算力。推测解码可以利用闲置的算力，分摊解码成本。
  \item \textbf{高接受率}。代码补全、结构化生成、事实型问答等任务的输出较容易预测，因而 \(\alpha\) 较高，加速幅度较大。
  \item \textbf{显存充足}。显存需要同时容纳草稿模型与目标模型。
\end{itemize}
\end{conceptbox}

\begin{notebox}{注意}
在批大小较大、GPU 已经饱和的服务场景中，推测解码\emph{无法提高吞吐量}。这项优化的目标是降低时延。
\end{notebox}

\begin{interview}
\question{1}{证明推测解码的输出与目标模型的输出具有相同的分布。}
证明依赖\textbf{拒绝采样}。设草稿模型以概率 \(q(x)\) 提出词元 \(x\)，目标模型的分布为 \(p(x)\)。以概率 \(\min(1,p(x)/q(x))\) 接受草稿词元。若拒绝，则从残余分布（residual distribution）\(p'(x)\propto\max(0,p(x)-q(x))\) 中采样。任一步的边际接受概率为
\[
  \sum_x q(x)\min(1,p(x)/q(x))=\sum_x\min(q(x),p(x))=1-\operatorname{TV}(p,q)\text{。}
\]
其中 \(\operatorname{TV}\) 表示总变差距离（total variation distance）。直接接受 \(x\) 的概率质量为 \(q(x)\min(1,p(x)/q(x))\)，校正步骤补足 \(p(x)\) 的剩余概率质量。两项合在一起，恰好得到分布 \(p(x)\)。

\question{2}{推测解码中的接受率 \(\alpha\) 是什么？哪些因素会影响它？}
接受率 \(\alpha\) 是单个草稿词元通过目标模型拒绝采样判据的概率，取决于草稿分布 \(q(x)\) 与目标分布 \(p(x)\) 的重叠程度，即 \(\alpha=1-\operatorname{TV}(p,q)=\sum_x\min(p(x),q(x))\)。提高接受率的因素包括以下几点。（1）\textbf{任务的可预测性}。编程和结构化文本的下一个词元更容易确定。
% <<<PDFPAGE 23>>>
（2）\textbf{草稿模型与目标模型的一致程度}。可选择同一系列的草稿模型，或由目标模型蒸馏得到的模型。（3）\textbf{采样温度}。温度较低时，两个分布都更集中，因而更可能一致。（4）\textbf{较小的推测窗口}。草稿词元数 \(K\) 较少时，逐个验证时遇到拒绝的累计概率也较低。在代码生成中，\(\alpha\) 通常为 \(0.7\) 至 \(0.9\)，开放式对话中通常为 \(0.5\) 至 \(0.7\)。

\question{3}{Medusa 与使用独立草稿模型的方法有何区别？两者各有什么取舍？}
\textbf{独立的草稿模型}是一个小模型，拥有自己的权重和 KV 缓存，并独立执行前向传播。它的优势是灵活。任何兼容的小模型都可以用作草稿模型，其质量也可以独立改进。代价是存储两个模型的显存开销，以及管理两份 KV 缓存的复杂性。\textbf{Medusa} 在目标模型的隐状态之上加上轻量线性头，无需独立模型。它直接使用目标模型的内部表示，这些表示比小型草稿模型的内部表示更丰富。另一方面，Medusa 头的表达能力弱于完整的草稿模型，因为头部只接收最后一个隐状态，无法直接访问完整上下文。因此，其接受率可能较低。Medusa-2 用自蒸馏缓解这个问题。Medusa 的部署更简单，精心选取的独立草稿模型则可能有更高的质量。

\question{4}{在什么情况下你不建议采用推测解码？}
我不建议在以下情形采用推测解码。（1）\textbf{批大小较大的生产服务}。批大小达到 32 或更大且 GPU 算力已经饱和时，并行验证会占用原本可用于处理更多请求的算力，导致吞吐量下降。（2）\textbf{受显存容量限制的环境}。无论是独立草稿模型还是 Medusa 头，都要占用额外显存。若显存已经全部分配给 KV 缓存和权重，就无法部署推测解码。（3）\textbf{接受率本来就低的任务}。开放式创作若采用高温度和多样的词汇，往往会导致较低的 \(\alpha\)（低于 \(0.5\)）。频繁拒绝会使加速幅度很小。（4）\textbf{输出极短的任务}。起草循环有固定开销。当输出只有 1 至 5 个词元时，推测解码没有收益。

\question{5}{描述 EAGLE 的架构，并解释它为何能达到比基本版 Medusa 更高的接受率。}
EAGLE 训练一个轻量的自回归草稿模型。它的输入既包含 Medusa 也会使用的目标模型隐状态 \(\mathbf{h}_t\)，又包含最近生成的词元的\textbf{词元嵌入} \(\mathbf{e}_t\)。将联合输入 \([\mathbf{h}_t;\mathbf{e}_t]\) 送入一个 Transformer 解码器层，得到下一步的隐状态 \(\mathbf{h}_{t+1}\) 的预测，再据此预测后续词元。相比 Medusa，它的主要优势是\textbf{自回归生成草稿}。EAGLE 可以依次生成多个草稿词元，每一步都以此前的词元为条件，无须从同一个固定隐状态出发预测全部词元。按顺序引入条件，能捕获 Medusa 的并行头未能利用的词元间依赖，使接受率高出 5 至 15 个百分点。EAGLE-2 还加入了基于置信度的动态树扩展。
\end{interview}
